\documentclass[11pt]{article}
\usepackage[margin=1in]{geometry}
\usepackage[T1]{fontenc}
\usepackage[utf8]{inputenc}
\usepackage{lmodern}
\usepackage{amsmath,amssymb,mathtools,bm}
\usepackage{booktabs}
\usepackage{enumitem}
\usepackage{hyperref}
\usepackage{microtype}

\title{Investigation plan for finite smooth thermal-expansion response in the non-Kramers pyrochlore model}
\author{}
\date{\today}

\begin{document}
\maketitle

\section{Purpose}

This note records what has been established numerically so far for the
16-site pyrochlore XXZ$+J_{3s}$ exact-diagonalization study, what symmetry
conclusions follow, and what controlled follow-up campaign is being launched
to test candidate explanations for a finite but slowly varying experimental
response.

\section{Established results so far}

\subsection{Model and observable}

The Hamiltonian studied numerically is the 16-site pyrochlore XXZ model with
an additional trilinear Kadowaki-type term:
\begin{equation}
H = H_{\rm XXZ} + H_{3s},
\end{equation}
with periodic boundary conditions and no external symmetry-breaking field.

The exact quantity computed is the connected quadrupolar energy covariance
\begin{equation}
\alpha_Q(T)
=
\frac{\langle QH\rangle_T - \langle Q\rangle_T \langle H\rangle_T}{T^2}
=
\frac{d\langle Q\rangle_T}{dT},
\end{equation}
for the five quadrupolar probes
\begin{equation}
Q_1,\; Q_2 \in E_g,
\qquad
Q_{xy},\; Q_{xz},\; Q_{yz} \in T_{2g}.
\end{equation}

\subsection{Exact ED findings}

The connected correlator was computed by full exact diagonalization in the full
$2^{16}=65536$ Hilbert space, using the exact $\mathbb Z_2\times\mathbb Z_2$
fcc translation decomposition into four momentum sectors and dense Hermitian
diagonalization in each sector.

The main numerical findings are:
\begin{enumerate}[leftmargin=1.8em]
\item For $J_{3s}=0$, all five $\alpha_Q(T)$ channels vanish to machine
precision. This is consistent with the internal pseudospin selection rule
of the XXZ limit.
\item For $J_{3s}=0.08$, the exact full-ED results still show only a tiny
residual signal in the $E_g$ channels and essentially zero signal in the
$T_{2g}$ channels:
\begin{equation}
\max_T |\alpha_{E_g}| \sim 10^{-8},
\qquad
\max_T |\alpha_{T_{2g}}| \sim 10^{-14},
\end{equation}
which is numerically negligible on the scale of the experimental response.
\item Therefore the exact cubic zero-field XXZ$+J_{3s}$ model does not provide
a meaningful finite non-scalar one-point thermal-expansion response.
\end{enumerate}

\subsection{Symmetry interpretation}

The exact results are consistent with the symmetry statement:
\begin{equation}
\text{cubic scalar Hamiltonian}
\quad \Longrightarrow \quad
\langle Q_{E_g}\rangle = \langle Q_{T_{2g}}\rangle = 0,
\end{equation}
and hence
\begin{equation}
\alpha_{E_g}(T)=\alpha_{T_{2g}}(T)=0,
\end{equation}
for a fully symmetry-preserving thermal trace.

The trilinear $J_{3s}$ term removes the accidental internal pseudospin
selection rule of the XXZ limit, but by itself it does not remove the cubic
point-group obstruction. Therefore a fully cubic scalar $J_{3s}$ term does not
generate a finite $E_g$ or $T_{2g}$ one-point response.

\section{Experimental motivation}

The experimental $[110]$ length-change signal mixes
\begin{equation}
A_{1g}\oplus E_g \oplus T_{2g}.
\end{equation}
The attached low-temperature dataset shows a finite, smooth, slowly varying
response with a very low characteristic temperature scale (visible structure
already near $T \sim 10^{-2}$ in the supplied temperature units), rather than a
clear symmetry-breaking divergence.

This points away from a clean spontaneous $E_g$ or $T_{2g}$ ordering transition
of the cubic XXZ$+J_{3s}$ Hamiltonian, and toward one of the following
possibilities:
\begin{enumerate}[leftmargin=1.8em]
\item a dominant $A_{1g}$ scalar background;
\item a weak explicit $E_g$ symmetry-breaking field inducing a smooth finite
$E_g$ response;
\item a weak explicit $T_{2g}$ symmetry-breaking field inducing a smooth finite
$T_{2g}$ response;
\item disorder acting as a random or biased non-cubic field;
\item a glassy or frozen quadrupolar response set by a low emergent disorder
scale rather than by the bare exchange scale.
\end{enumerate}

\section{Working interpretation}

The clean-limit exact ED strongly suggests that any realistic finite non-scalar
experimental response must come from \emph{induced} rather than spontaneous
one-point polarization:
\begin{equation}
H = H_{A_{1g}} + \lambda_\Gamma V_\Gamma,
\qquad
\Gamma = E_g \text{ or } T_{2g}.
\end{equation}
Then, to leading order,
\begin{equation}
\langle Q_\Gamma \rangle \sim \lambda_\Gamma \chi_\Gamma(T),
\qquad
\alpha_\Gamma(T) \sim \lambda_\Gamma \frac{d\chi_\Gamma(T)}{dT}.
\end{equation}
If $\lambda_\Gamma$ is small and $\chi_\Gamma(T)$ is smooth and noncritical,
the resulting response is finite but slowly varying, precisely the regime that
appears experimentally plausible.

\section{Follow-up campaign: what to test}

\subsection{Main campaign objective}

The immediate goal is to determine which small explicit symmetry-breaking source
can generate a finite, smooth low-energy response in the exact 16-site model.

\subsection{Campaign hypotheses}

The campaign will test the following hypotheses.

\paragraph{H1: weak explicit $E_g$ source}
Add
\begin{equation}
H \rightarrow H - \lambda_E Q_1
\quad\text{or}\quad
H \rightarrow H - \lambda_E Q_2.
\end{equation}
Prediction: finite induced $E_g$ response with $T_{2g}$ remaining zero at
linear order.

\paragraph{H2: weak explicit $T_{2g}$ source}
Add
\begin{equation}
H \rightarrow H - \lambda_T Q_{xy}
\quad\text{or}\quad
H \rightarrow H - \lambda_T Q_{xz}
\quad\text{or}\quad
H \rightarrow H - \lambda_T Q_{yz}.
\end{equation}
Prediction: finite induced $T_{2g}$ response with $E_g$ remaining zero at
linear order.

\paragraph{H3: disorder / random-field origin}
This is not yet launched in the first pilot, but the source-field pilot is the
clean prerequisite for it. If neither uniform $E_g$ nor uniform $T_{2g}$ bias
gives a plausible smooth signal, the next stage should introduce random local
transverse fields or random local $E_g/T_{2g}$ strain fields.

\subsection{Pilot grid being launched}

The first-pass pilot uses the exact full-ED source-field Hamiltonian with
\begin{equation}
\lambda = 10^{-3}
\end{equation}
for the following four cases:
\begin{align}
\texttt{EgQ1\_1e-3}:& \quad H \rightarrow H - 10^{-3} Q_1, \\
\texttt{EgQ2\_1e-3}:& \quad H \rightarrow H - 10^{-3} Q_2, \\
\texttt{T2gXY\_1e-3}:& \quad H \rightarrow H - 10^{-3} Q_{xy}, \\
\texttt{T2gXZ\_1e-3}:& \quad H \rightarrow H - 10^{-3} Q_{xz}.
\end{align}

This pilot is chosen to answer the most immediate question with minimal extra
wall-clock cost:
\begin{quote}
Can a very small explicit $E_g$ or $T_{2g}$ bias produce a finite but smooth
induced thermal-expansion response of the kind seen experimentally?
\end{quote}

\section{Numerical implementation}

The exact ED driver has been extended to accept repeated source terms of the
form
\begin{equation}
\texttt{--source NAME=VALUE},
\end{equation}
which apply the modification
\begin{equation}
H \leftarrow H - \texttt{VALUE}\times Q_{\texttt{NAME}}.
\end{equation}

The campaign is launched by the shell script
\begin{equation}
\texttt{twist\_qsi\_demo/scripts/launch\_source\_field\_campaign.sh},
\end{equation}
which runs the pilot source-field cases under \texttt{nohup} with
checkpointing and separate logs.

\section{Planned analysis after launch}

For each source-field case, the following comparisons will be made:
\begin{enumerate}[leftmargin=1.8em]
\item magnitude and sign of the induced $\langle Q\rangle(T)$;
\item induced $\alpha_Q(T)$ in the directly sourced channel;
\item cross-talk into the other irreps;
\item smoothness of the low-$T$ profile;
\item whether the characteristic temperature scale is suppressed relative to
the bare exchange scale.
\end{enumerate}

If the pilot identifies a plausible induced-response scenario, the next grid
should vary $\lambda$ over
\begin{equation}
10^{-4},\; 3\times 10^{-4},\; 10^{-3},\; 3\times 10^{-3},\; 10^{-2},
\end{equation}
to test linear-response scaling and to compare against the experimental
amplitude and smoothness.

If the pilot fails to reproduce the qualitative response, the next campaign
should move to disorder realizations and sample-to-sample distributions rather
than uniform source fields.

\section{Bottom line}

The clean cubic XXZ$+J_{3s}$ model does not produce a meaningful finite
$E_g/T_{2g}$ one-point thermal-expansion response in exact zero-field ED.
Therefore the experimentally observed smooth finite low-energy signal is most
naturally interpreted as either
\begin{enumerate}[leftmargin=1.8em]
\item a dominant scalar $A_{1g}$ background, or
\item a weak induced non-scalar response from explicit symmetry breaking or
disorder.
\end{enumerate}
The launched source-field campaign is the first controlled numerical test of
that induced-response scenario.

\end{document}
